\section{Residual mesh parameters and deformation of module categories}
\label{sec:positive-deformations}

From now on we assume that $A$ is basic and $\charac k=2$, retaining the
graded category $\C=\gr_\JJ\ind(\modu A)$ and its standard form $A_0$ from
Part~\ref{part:graded-form}. In characteristic two, the mesh derivative
$D_m(\C)$ need not be surjective, so higher-degree error terms in the mesh
relations cannot in general be eliminated by adjusting the irreducible
morphisms. We begin by describing the cokernel of $D_m(\C)$ on the
$\tau$-periodic components of $\Gamma_A$. We then introduce positive
filtered deformations of $A_0$—of which $A$ itself is the primary
example—and establish internal-degree bounds on $\Ext$ groups that allow us
to lift graded $A_0$-modules and their morphisms to any positive filtered
deformation. In Section~\ref{sec:weights-and-parameters} we use this module
lifting to embed the positive-degree Hochschild cohomology of $A_0$ into the
cokernel of the mesh derivative.

\subsection{The remaining linear correction spaces}

By Proposition~\ref{prop:periodic-localization}, the cokernel of the mesh
derivative $D_m(\C)$ of \eqref{eq:mesh-derivative} is isomorphic to
$\operatorname{coker}D_m(\C_T)$ on the stable subquiver of $\tau$-periodic
vertices. Evaluating $\operatorname{coker}D_m$ on the periodic tree models
of Proposition~\ref{prop:periodic-tree-model} in characteristic two yields
the following classification, proved in
Appendix~\ref{app:characteristic-two-calculations} (where periodic
components of type $A$ are shown to have vanishing cokernel in
Appendix~\ref{app:char2-line-models}).

\begin{proposition}
\label{prop:residual-correction-spaces}
Let $C=C(Q,c)$ be a finite-dimensional scalar mesh category on a finite
translation quiver $Q$ with at most one arrow between each ordered pair of
vertices. Assume that $\charac k=2$ and that $C$ satisfies the brick
hypothesis of Subsection~\ref{subsec:brick-hypothesis}. Call a periodic
component of $Q$ containing at least one arrow \emph{residual} if
$\bigoplus_{m\geq2}\operatorname{coker}D_m(C)$ has a nonzero class supported
on that component, and let $I$ index the residual components of $Q$.
\begin{enumerate}[label={\upshape(\arabic*)}]
\item The $i$th residual component ($i\in I$) is isomorphic to
$\Z\Delta/\langle\tau^{s_i}\rho\rangle$, where $\Delta$ is a Dynkin tree of
type $D$ or $E$ and $s_i\geq1$, as in
Proposition~\ref{prop:periodic-tree-model}(2).
\item For each $i\in I$, there exist a vertex $v_i$ of the $i$th residual
component and a path $p_i:\tau v_i\to v_i$ of length $d_i=2s_i+2$ such that,
viewing $p_i$ as an element of $V_{d_i}(C)$ supported in the row of $v_i$
and writing $[p_i]$ for its class in $\operatorname{coker}D_{d_i-1}(C)$, we
have
\[
 \bigoplus_{m\geq2}\operatorname{coker}D_m(C)
                      =\bigoplus_{i\in I}k[p_i].
\]
In particular, $d_i=2s_i+2$ is the unique path degree in which the $i$th
residual component has a nonzero cokernel class.
\end{enumerate}
\end{proposition}

By Subsection~\ref{subsec:brick-hypothesis}, Proposition~\ref{prop:residual-correction-spaces}
applies directly to $C=\C$.

\subsection{Degrees of extensions between normalized modules}
\label{subsec:ext-degrees}

To relate filtered deformations of $A_0$ to the mesh cokernel classes
$[p_i]$, we will lift the graded indecomposable $A_0$-modules to any
positive filtered deformation $B$ of $A_0$.

\begin{definition}
\label{def:positive-deformation}
A \emph{positive filtered deformation} of a finite-dimensional nonnegatively
graded algebra $\Lambda=\bigoplus_{d\geq0}\Lambda_d$ is an algebra $B$
equipped with a finite decreasing multiplicative filtration $F^\bullet B$
and a specified graded algebra isomorphism $\gr_F B\cong\Lambda$.
\end{definition}

In particular, the categorical radical filtration
$F^n A(x,y)=\JJ^n(P_y,P_x)$ of Section~\ref{sec:projective-rays} makes the
original algebra $A$ a positive filtered deformation of its standard form
$A_0$.

In classical deformation theory \cite{Gerstenhaber,Lowen}, the obstructions
to lifting a graded $\Lambda$-module action to a positive filtered
deformation lie in $\Ext^2_\Lambda$ in positive internal degrees, while the
choices of module lifts and the obstructions to lifting morphisms lie in
$\Ext^1_\Lambda$ in positive internal degrees. Here, for graded right
$\Lambda$-modules $M$ and $N$, a homomorphism or extension has
\emph{internal degree} $t\in\Z$ if it is represented on a graded projective
resolution of $M$ by a cochain that shifts degrees by $+t$. Using the shift
convention $(N\langle a\rangle)_q=N_{q-a}$ (so that a degree-zero map
$M\langle a\rangle\to N\langle b\rangle$ has underlying internal degree
$a-b$), we have
\[
 \Ext^n_\Lambda(M,N)_t
   =\Ext^n_{\mathrm{gr}\text{-}\Lambda}(M,N\langle-t\rangle),
\]
where $\Ext_{\mathrm{gr}\text{-}\Lambda}$ is computed in the category of
graded right $\Lambda$-modules with degree-zero morphisms.

For $\Lambda=A_0$, enumerate the objects of $\C$ as $X_1,\ldots,X_h$, and
let $M_1,\ldots,M_h$ be the corresponding graded indecomposable right
$A_0$-modules obtained by projective evaluation, graded by
$(M_i)_d e_x=\C_d(P_x,X_i)$. We refer to $M_1,\ldots,M_h$ as the
\emph{degree-normalized} indecomposable $A_0$-modules; note that the
indecomposable projective $A_0$-modules $e_x A_0$ carry their natural
gradings under this normalization. Full faithfulness of projective
evaluation (Proposition~\ref{prop:projective-evaluation}(1)) immediately
gives
\begin{equation}
 \Hom_{A_0}(M_i,M_j)_r=0\quad(r<0),\qquad
 \Hom_{A_0}(M_i,M_j)_0=
 \begin{cases}k\cdot\id_{M_i},&\text{if }i=j,\\0,&\text{if }i\ne j.\end{cases}
 \label{eq:normalized-hom-degrees}
\end{equation}

\begin{lemma}
\label{lem:positive-ext-vanishing}
Let $A_0$ be the standard form of Construction~\ref{con:projective-algebra},
and let $M_1,\ldots,M_h$ be the degree-normalized indecomposable
$A_0$-modules defined above.
\begin{enumerate}[label={\upshape(\arabic*)}]
\item Every finite-dimensional graded right $A_0$-module is isomorphic to a
direct sum of degree shifts of the modules $M_i$.
\item For all $i,j$ and all $n\geq1$,
\[
                 \Ext^n_{A_0}(M_i,M_j)_t=0\qquad\text{for all }t>-2n.
\]
\end{enumerate}
\end{lemma}
\begin{proof}
Let $X$ be a finite-dimensional graded right $A_0$-module that is
indecomposable as a graded module. Decomposing $X$ as an ungraded
$A_0$-module into copies of the indecomposable modules $M_i$, let
$p_a:X\to M_{i_a}$ and $j_a:M_{i_a}\to X$ be the projection and inclusion
maps of the summands, so that $1_X=\sum_a p_a j_a$ in path order. Passing to
the homogeneous component of degree zero gives
\[
                   1_X=\sum_{a,r}(p_a)_r(j_a)_{-r}
\]
in the local ring $\End_{\mathrm{gr}\text{-}A_0}(X)$. Some single term
$(p_a)_r(j_a)_{-r}$ is therefore invertible, exhibiting $X$ as a graded
direct summand of $M_{i_a}\langle-r\rangle$. Since $M_{i_a}$ is already
indecomposable as an ungraded module, $X\cong M_{i_a}\langle-r\rangle$,
proving (1).

We prove (2) first for $n=1$. The vanishing is trivial if $M_i$ is
projective. If $X_i$ is nonprojective, Proposition~\ref{prop:graded-realization}(1)
and Lemma~\ref{lem:mesh-simple-tops} provide an exact sequence of
representable functors
\[
 0\longrightarrow\C(-,\tau X_i)\longrightarrow
 \bigoplus_{Y\to X_i}\C(-,Y)\longrightarrow\C(-,X_i)
\]
whose cokernel is the simple functor $\C_0(-,X_i)$. Restricting these
functors to the projective objects $P_x$ annihilates the simple cokernel
(since $X_i$ is nonprojective) and yields a short exact sequence of graded
$A_0$-modules with degree-one maps. Under the equivalence
$\add\C\simeq\modu A_0$, this short exact sequence is the almost split
sequence ending at $M_i$; inserting degree shifts so that its maps become
homogeneous of degree zero writes it as
\[
 0\longrightarrow M_{\tau i}\langle2\rangle\longrightarrow
 \bigoplus_{Y\to X_i}M_Y\langle1\rangle\longrightarrow M_i
 \longrightarrow0.
\]
Now suppose for a contradiction that
$0\to M_j\langle-t\rangle\to E\xrightarrow{v}M_i\to0$ is a nonsplit graded
extension with $t>-2$. It remains nonsplit as an ungraded sequence (since any
ungraded section would have a degree-zero component that is a graded
section). By the lifting property of the right almost split map, $v$ factors
through $\bigoplus_{Y\to X_i}M_Y\langle1\rangle\to M_i$, and taking the
degree-zero component produces a graded factorization. The induced graded
map on kernels $M_j\langle-t\rangle\to M_{\tau i}\langle2\rangle$ has
underlying internal degree $-t-2<0$ and therefore vanishes by
\eqref{eq:normalized-hom-degrees}. Consequently, the factorization descends
to a section of the almost split sequence ending at $M_i$, a contradiction.
Thus $\Ext^1_{A_0}(M_i,M_j)_t=0$ for all $t>-2$.

For $n\geq2$, observe that in a minimal graded projective presentation of a
nonprojective module $M_i$, every projective summand $P\langle a\rangle$ of
the projective cover has $a\geq1$ because its nonzero map to $M_i$ lies in
positive degree; similarly, every indecomposable summand $Z\langle b\rangle$
of the first syzygy $\Omega^1 M_i$ maps into $P\langle a\rangle$ via a
radical morphism of positive degree, so $b-a\geq1$ and hence $b\geq2$. By
induction on $r$, every indecomposable summand of the $r$th syzygy
$\Omega^r M_i$ has the form $Z\langle b\rangle$ with $b\geq 2r$. Dimension
shifting then gives
\[
 \Ext^n_{A_0}(M_i,M_j)_t
 =\Ext^1_{A_0}(\Omega^{n-1}M_i,M_j)_t=0
 \qquad\text{for all }t>-2n,
\]
because each indecomposable summand $Z\langle b\rangle$ of
$\Omega^{n-1}M_i$ has $b\geq 2(n-1)$ and thus contributes
$\Ext^1_{A_0}(Z,M_j)_{t+b}=0$ since $t+b>-2$.
\end{proof}

\subsection{Lifting actions and morphisms}

Before lifting modules to a positive filtered deformation $B$, we place $B$
in a split relative normal form over its degree-zero semisimple subalgebra.

For an algebra $\Lambda$ over a commutative ring $R$ and a unital subalgebra
$S\subseteq\Lambda$, the space of \emph{normalized $S$-relative Hochschild
$n$-cochains} ($n\geq1$) is
\begin{equation}
 C^n_S(\Lambda,\Lambda)
   =\Hom_{S\text{-}S}\bigl((\Lambda/S)^{\otimes_S n},\Lambda\bigr),
 \label{eq:normalized-relative-cochains}
\end{equation}
with $C^0_S(\Lambda,\Lambda)=\{a\in\Lambda:as=sa\text{ for all }s\in S\}$
in degree zero and the standard Hochschild differential (in particular,
$(df)(a,b)=af(b)-f(ab)+f(a)b$ for a $1$-cochain $f$). Viewed as an
$R$-multilinear map on $\Lambda^n$, a normalized $S$-relative $n$-cochain is
$S$-bilinear in the outer variables, $S$-balanced across each adjacent pair
of inputs, and vanishes whenever any input belongs to $S$. We will use these
cochains over fields and over $\Z/4$.

\begin{lemma}
\label{lem:split-positive-deformation}
Let $k$ be a field, let $\Lambda=\bigoplus_{d\geq0}\Lambda_d$ be a
finite-dimensional nonnegatively graded $k$-algebra whose degree-zero part
$S=\Lambda_0$ is split semisimple, and let $B$ be a positive filtered
deformation of $\Lambda$. Then $S$ lifts to a unital subalgebra of $B$, and
the filtration $F^\bullet B$ splits as a filtration of $S$-bimodules. Under
the resulting filtered $S$-bimodule identification $B\cong\Lambda$ (inducing
the given identification $\gr_F B=\Lambda$), the multiplication of $B$ takes
the form $\mu_0+\sum_{q>0}\mu_q$, where $\mu_0$ is the multiplication of
$\Lambda$ and each $\mu_q\in C^2_S(\Lambda,\Lambda)$ is a normalized
$S$-relative $2$-cochain raising internal degree by $q$.
\end{lemma}
\begin{proof}
Set $\mathfrak n=F^1 B$, so that $\mathfrak n$ is a nilpotent two-sided
ideal of $B$ with $B/\mathfrak n\cong S$. Because idempotents can be lifted
modulo any nilpotent ideal—iterating $a\mapsto a+(1-2a)(a^2-a)$ across
successive square-zero quotients—and orthogonal idempotents lift
successively in complementary corners, we may lift all diagonal matrix units
of the matrix blocks of $S$ to a complete set of orthogonal idempotents in
$B$.

Within a single matrix block $M_n(k)$ of $S$, let $p_1,\ldots,p_n\in B$ be
the lifted diagonal idempotents, and choose elements $u_i\in p_i B p_1$ and
$v_i\in p_1 B p_i$ lifting the elementary matrix units $E_{i1}$ and $E_{1i}$
(with $u_1=v_1=p_1$). Since $v_i u_i\equiv p_1\pmod{\mathfrak n}$, the
element $v_i u_i$ is invertible in $p_1 B p_1$; replacing $v_i$ by
$(v_i u_i)^{-1}v_i$ ensures that $v_i u_i=p_1$. Then $u_i v_i$ is an
idempotent in $p_i B p_i$ congruent to $p_i$ modulo $\mathfrak n$, so
$p_i-u_i v_i$ is an idempotent in $\mathfrak n$ and therefore vanishes. Thus
$u_i v_i=p_i$, and the products $u_i v_j$ satisfy the matrix-unit relations
of $M_n(k)$. Since idempotents in distinct matrix blocks are orthogonal,
these matrix units define a unital algebra embedding $S\hookrightarrow B$
lifting $B/\mathfrak n\cong S$.

Finally, because $S\otimes_k S^{\op}$ is a finite product of matrix algebras
over $k$, it is semisimple, so every short exact sequence of $S$-bimodules
splits. Choosing $S$-bimodule splittings $F^d B/F^{d+1}B\hookrightarrow F^d B$
(using the subalgebra $S$ itself in degree $d=0$) identifies $B$ with
$\Lambda$ as a filtered $S$-bimodule inducing the prescribed isomorphism
$\gr_F B\cong\Lambda$. Under this identification, the multiplication of $B$
is $S$-bilinear, restricts to the multiplication of $S$ on $S$, and raises
grading degree by at least $0$, with degree-zero component equal to $\mu_0$.
\end{proof}

\begin{proposition}
\label{prop:positive-module-lifting}
Let $k$ be a field, let $\Lambda$ be a finite-dimensional nonnegatively
graded $k$-algebra with split semisimple degree-zero part $S=\Lambda_0$, and
let $M_1,\ldots,M_h$ be finite-dimensional graded right $\Lambda$-modules
satisfying
\[
 \Ext^2_\Lambda(M_i,M_i)_q=0,\qquad
 \Ext^1_\Lambda(M_i,M_j)_q=0\qquad\text{for all }i,j\text{ and all }q>0.
\]
Let $B$ be a positive filtered deformation of $\Lambda$.
\begin{enumerate}[label={\upshape(\arabic*)}]
\item Each $M_i$ admits a filtered right $B$-module lift $M_i^B$ with
$\gr M_i^B=M_i$, unique up to a filtered $B$-module isomorphism inducing the
identity on $M_i$.
\item Every homogeneous $\Lambda$-module homomorphism $M_i\to M_j$ of degree
$r\geq0$ lifts to a filtered $B$-module homomorphism $M_i^B\to M_j^B$ in
filtration degree $r$ with that prescribed leading term. Consequently, the
associated graded Hom spaces in nonnegative degrees between the lifted
modules coincide with the original graded Hom spaces.
\item Suppose further that $\Hom_\Lambda(M_i,M_j)_r=0$ for all $r<0$ and
$\Hom_\Lambda(M_i,M_j)_0=\delta_{ij}k$. Then the full subcategory $E$ of
$\modu B$ on $M_1^B,\ldots,M_h^B$, filtered by map degree, has associated
graded category equal to the graded Hom category on $M_1,\ldots,M_h$; if the
latter is generated in degree one, then the map-degree filtration on $E$
coincides with its categorical radical filtration. Moreover, if
$M_i=e_x\Lambda$ for a primitive idempotent $e_x\in S$, we may take
$M_i^B=e_x B$.
\end{enumerate}
\end{proposition}
\begin{proof}
Identify $B$ as a filtered $S$-bimodule with $\Lambda$ via
Lemma~\ref{lem:split-positive-deformation}, so that the multiplication $a*b$
of $B$ has the form $\mu_0(a,b)+\sum_{q>0}\mu_q(a,b)$. Because $S$ is
semisimple, any filtered right $B$-module lift of $M_i$ can likewise be
identified with $M_i$ as a filtered right $S$-module, so that its $B$-action
takes the form $m\cdot a=\nu_0(m,a)+\sum_{q>0}\nu_q(m,a)$, where $\nu_0$ is
the $\Lambda$-action on $M_i$ and each $\nu_q:M_i\otimes_S(\Lambda/S)\to M_i$
is a right $S$-linear map raising internal degree by $q$. Recall that the
normalized $S$-relative bar resolution of a graded right $\Lambda$-module $M$
has degree-$n$ term
\[
           M\otimes_S(\Lambda/S)^{\otimes_S n}\otimes_S \Lambda,
\]
which is a graded projective right $\Lambda$-module because $S$ is
semisimple; applying $\Hom_\Lambda(-,N)$ computes $\Ext^n_\Lambda(M,N)$ via
normalized $S$-relative bar cochains.

Suppose by induction on $t\geq1$ that we have chosen $\nu_1,\ldots,\nu_{t-1}$
so that the associativity identity $(m\cdot a)\cdot b=m\cdot(a*b)$ holds in
internal degrees less than $t$. A direct check using the associativity of
$*$ and the lower-degree action identities shows that the degree-$t$ failure
of $(m\cdot a)\cdot b=m\cdot(a*b)$ is a normalized $S$-relative bar
$2$-cocycle of internal degree $t$ with values in $\Hom_k(M_i,M_i)$. Since
$\Ext^2_\Lambda(M_i,M_i)_t=0$, this $2$-cocycle is the bar coboundary of a
normalized $1$-cochain $\nu_t$ of internal degree $t$, which cancels the
degree-$t$ error. Because $M_i$ is finite-dimensional, this induction
terminates after finitely many steps and yields a filtered $B$-module action
on $M_i$. Similarly, if two such actions agree in degrees less than $t$,
their degree-$t$ difference is a bar $1$-cocycle of internal degree $t$; by
$\Ext^1_\Lambda(M_i,M_i)_t=0$, it is the coboundary of an $S$-linear
endomorphism $h_t$ of $M_i$ of degree $t$, so conjugating by the invertible
filtered automorphism $1+h_t$ brings the two actions into agreement through
degree $t$. This proves (1).

For (2), let $f_0:M_i\to M_j$ be homogeneous of degree $r\geq0$, and suppose
we have constructed $f=f_0+f_1+\cdots+f_{t-1}$ (with $f_s$ of degree $r+s$)
commuting with the $B$-actions in degrees less than $r+t$. The degree-$(r+t)$
error in $f(m\cdot a)=f(m)\cdot a$ is a bar $1$-cocycle in
$C^1_S(M_i,M_j)$ of internal degree $r+t>0$. Since
$\Ext^1_\Lambda(M_i,M_j)_{r+t}=0$, this error is the coboundary of a map
$f_t:M_i\to M_j$ of degree $r+t$, allowing us to continue the induction to
obtain a filtered $B$-linear lift of $f_0$. Since the leading homogeneous
term of any nonzero filtered $B$-module homomorphism $M_i^B\to M_j^B$ is a
nonzero homogeneous $\Lambda$-module homomorphism $M_i\to M_j$, this proves
(2).

Under the hypotheses of (3), there are no nonzero homomorphisms of negative
degree, so the associated graded category of $E$ is the full graded Hom
category on $M_1,\ldots,M_h$. The positive-degree filtration ideal $F^1 E$
is nilpotent and $E/F^1 E\cong\prod_i k$, so $F^1 E$ is the categorical
radical $\JJ(E)$ of $E$ (and in particular the lifted modules $M_i^B$ are
indecomposable and pairwise nonisomorphic). If $\gr_F E$ is generated in
degree one, then each $F^q E/F^{q+1}E$ is spanned by products of $q$
elements of $F^1 E/F^2 E$, so descending induction on $q$ gives
$F^q E=\JJ(E)^q$ for all $q\geq1$. Finally, when $M_i=e_x\Lambda$, the lifted
projective module $e_x B$ has associated graded module $e_x\Lambda$, so by
the uniqueness in (1) we may take $M_i^B=e_x B$.
\end{proof}

Applying Proposition~\ref{prop:positive-module-lifting} to $\Lambda=A_0$ and
its degree-normalized indecomposable modules $M_1,\ldots,M_h$ (whose Hom and
$\Ext$ groups satisfy the required bounds by \eqref{eq:normalized-hom-degrees}
and Lemma~\ref{lem:positive-ext-vanishing}(2)) shows that for every positive
filtered deformation $B$ of $A_0$, the lifted modules $M_i^B$ form a
filtered category with associated graded category $\C$, and its projective
corner recovers $\End_B(\bigoplus_x e_x B)^{\op}=B$.

In Section~\ref{sec:weights-and-parameters} and
Section~\ref{sec:integral-deformations} we also need the first-order version
of this lifting calculation over $R[\varepsilon]/(\varepsilon^2)$, where the
coefficient ring $R$ is either a field or $\Z/4$.

\begin{lemma}
\label{lem:square-zero-selected-lifts}
Let $R$ be a field or $\Z/4$, let $\Lambda$ be a finitely generated free
nonnegatively graded $R$-algebra with $\Lambda_0=S=R^n$, and let
$M_1,\ldots,M_h$ be finitely generated graded right $\Lambda$-modules that
are free over $R$. Suppose that each graded module $\Hom_\Lambda(M_i,M_j)$ is
free over $R$ and concentrated in nonnegative degrees, and that
\[
 \Ext^2_\Lambda(M_i,M_i)_t=0,\qquad
 \Ext^1_\Lambda(M_i,M_j)_t=0\qquad\text{for all }i,j\text{ and all }t>0.
\]
Let $\mu\in C^2_S(\Lambda,\Lambda)$ be a normalized $S$-relative Hochschild
$2$-cocycle of internal degree $q>0$, and form the graded first-order
deformation
$\Lambda_\mu=(\Lambda\otimes_R R[\varepsilon]/(\varepsilon^2),m+\varepsilon\mu)$
with $\deg\varepsilon=-q$. Then:
\begin{enumerate}[label={\upshape(\arabic*)}]
\item Each $M_i$ lifts to a graded right $\Lambda_\mu$-module $M_i^\mu$ free
over $R[\varepsilon]/(\varepsilon^2)$ (taking $M_i^\mu=e_x\Lambda_\mu$ when
$M_i=e_x\Lambda$), uniquely up to a degree-zero $\Lambda_\mu$-module
isomorphism reducing to the identity modulo $\varepsilon$.
\item Every homogeneous $\Lambda$-linear map $f:M_i\to M_j$ lifts to a
homogeneous $\Lambda_\mu$-linear map $M_i^\mu\to M_j^\mu$ of the same degree,
and for any homogeneous $R$-basis of $\Hom_\Lambda(M_i,M_j)$, any chosen
homogeneous lifts form an $R[\varepsilon]/(\varepsilon^2)$-basis of
$\Hom_{\Lambda_\mu}(M_i^\mu,M_j^\mu)$.
\item If $\Lambda$ and the $M_i$ carry a compatible grading by an abelian
group $L$ and $\mu$ is homogeneous of $L$-degree $\chi$, then assigning
$\varepsilon$ the $L$-degree $-\chi$ allows all module and morphism lifts to
be chosen homogeneous for the $L$-grading as well.
\end{enumerate}
\end{lemma}
\begin{proof}
Because the vertex components of $\Lambda$ and $M_i$ are direct summands of
finitely generated free $R$-modules, they are projective over $R$, so the
normalized $S$-relative bar resolution of $M_i$ consists of graded projective
$\Lambda$-modules and computes $\Ext^\bullet_\Lambda(M_i,M_j)$ over $R$ (including
when $R=\Z/4$). Writing a candidate lifted action on
$M_i\otimes_R R[\varepsilon]/(\varepsilon^2)$ as
$m\cdot a=ma+\varepsilon\nu_i(m,a)$ with $\nu_i\in C^1_S(M_i,M_i)_q$, the
associativity condition modulo $\varepsilon^2$ is
\[
 \nu_i(ma,b)-\nu_i(m,ab)+\nu_i(m,a)b=m\mu(a,b).
\]
Since $d\mu=0$, the right-hand side is a relative bar $2$-cocycle of internal
degree $q>0$, so the vanishing of $\Ext^2_\Lambda(M_i,M_i)_q$ guarantees the
existence of $\nu_i$. The difference of any two solutions is a relative bar
$1$-cocycle of degree $q$, which by $\Ext^1_\Lambda(M_i,M_i)_q=0$ is the
coboundary of a degree-$q$ endomorphism $h$ of $M_i$, so the two lifts are
identified by the automorphism $1+\varepsilon h$. This proves (1).

Similarly, a homogeneous map $f:M_i\to M_j$ of degree $r\geq0$ lifts to a
$\Lambda_\mu$-linear map $f+\varepsilon g$ if and only if
\[
 g(ma)-g(m)a=\nu_j(f(m),a)-f(\nu_i(m,a)).
\]
The right-hand side is a relative bar $1$-cocycle of internal degree
$r+q>0$, so the vanishing of $\Ext^1_\Lambda(M_i,M_j)_{r+q}$ provides the correction $g$.
If $F_1,\ldots,F_s$ are homogeneous lifts of a homogeneous $R$-basis
$f_1,\ldots,f_s$ of the entire module $\Hom_\Lambda(M_i,M_j)$, then
subtracting from any $\Phi\in\Hom_{\Lambda_\mu}(M_i^\mu,M_j^\mu)$ the $R$-linear
combination of $F_1,\ldots,F_s$ matching its reduction modulo $\varepsilon$
leaves an element of the form $\varepsilon H$ where $H\in\Hom_\Lambda(M_i,M_j)$
(when $\Phi$ is homogeneous of degree $r$, its reduction modulo $\varepsilon$
lies in degree $r$ and $H$ lies in degree $r+q$), which is therefore an
$R$-linear combination of $\varepsilon F_1,\ldots,\varepsilon F_s$. Conversely, if
$\sum_l(a_l+\varepsilon b_l)F_l=0$ with $a_l,b_l\in R$, reducing modulo
$\varepsilon$ gives $\sum_l a_l f_l=0$, so $a_l=0$; the remaining equation
$\varepsilon\sum_l b_l f_l=0$ in the free $R[\varepsilon]/(\varepsilon^2)$-module
$\Hom_R(M_i,M_j)\otimes_R R[\varepsilon]/(\varepsilon^2)$ then forces
$\sum_l b_l f_l=0$ and hence $b_l=0$. This proves (2). Finally, since the
bar differentials preserve the $L$-grading, projecting any solution onto its
homogeneous component of the required $L$-degree gives (3).
\end{proof}
